\chapter*{补充材料阅读地图}
\addcontentsline{toc}{chapter}{补充材料阅读地图}

本册主线只保留 12 个编号章。补充材料不再夹在第二章或第三章后面，否则读者刚建立的因果线会被大块细节打断。下面这些材料保留在源码树中，用作选读、改写素材和实践卷素材。

\begin{longtable}{@{}L{0.22\linewidth}L{0.28\linewidth}L{0.42\linewidth}@{}}
\toprule
材料类型 & 放置位置 & 使用方式 \\
\midrule
硬件深讲 & \filepath{source/latex/supplements/hardware/} & 作为计算机组成和硬件补充卷素材。需要追 ALU、最小 CPU、指令编码、ABI、cache/TLB、总线、DMA、IOMMU、APIC、PCIe 和平台设备细节时再读。主线第二章只引用结论，不连续展开这些文件。 \\
裸机前置 & \filepath{source/latex/supplements/pre-os/} & 作为“从裸机到内核入口”的选读素材。适合在读完第二章之后补 MCU 主循环、中断、定时器、启动、链接脚本和早期初始化，但不再作为 OS 主线的隐形章节。 \\
可运行模型与验收 & \filepath{source/latex/supplements/models/} & 作为实践与代码卷素材。这里保存 C++20 硬件模型、x86-64 契约模型、第一阶段验收和硬件到 OS 对象映射。主线正文只保留要验证的不变量，完整模型放到实践卷展开。 \\
源码走读 & \filepath{source/latex/chapters/} 中未直接接入主线的 deep/source/lab 文件 & 调度、同步、系统调用、page fault、内存、块层、文件系统、网络、恢复、观测和参考实现的长篇材料保留为选读与改写素材。主线各章已经压缩吸收其核心账本，不再把这些文件连续插入目录。 \\
\bottomrule
\end{longtable}

\begin{rubricbox}
补充材料的使用原则：先读主线，能说清“失败场景、硬件事实、内核对象、状态转换、失败回滚和证据”之后，再按问题查补充材料。不要把补充材料当成第二章后面的连续正文。
\end{rubricbox}

硬件深讲建议顺序如下：

\begin{enumerate}
\tightlist
\item \filepath{ch00-hardware-contract-map.tex}：先建立读硬件的统一方法。
\item \filepath{ch00-hardware-gates-to-cpu.tex}、\filepath{ch00-digital-logic-to-cpu-deep.tex}、\filepath{ch00-minimal-computer-from-scratch.tex}：从电路、寄存器和时钟走到最小 CPU。
\item \filepath{ch00-number-code-assembly-deep.tex}、\filepath{ch00-call-stack-abi-deep.tex}、\filepath{ch00-privileged-machine-walkthrough.tex}：把编号、汇编、调用栈、ABI 和特权入口接起来。
\item \filepath{ch00-cache-tlb-memory-deep.tex}、\filepath{ch00-memory-bus-io-deep.tex}、\filepath{ch00-microarchitecture-platform-deep.tex}、\filepath{ch00-platform-devices-timing-deep.tex}：再进入现代平台、缓存、地址翻译、总线、设备和时钟。
\item \filepath{ch00-hardware-os-casebook.tex}、\filepath{ch00-hardware-os-walkthroughs-stage-one.tex}、\filepath{ch00-hardware-case-lab.tex}：最后用端到端案例检查理解。
\end{enumerate}

裸机前置建议顺序如下：

\begin{enumerate}
\tightlist
\item \filepath{ch02-bare-metal-mcu.tex}：从温控板主循环看为什么轮询和阻塞函数会失败。
\item \filepath{ch03-interrupts-timers.tex}：从串口丢字节和坏任务死循环推导中断、定时器和抢占入口。
\item \filepath{ch04-boot-linker-kernel-init.tex}：从上电后不能直接调用 C 函数，推导 reset handler、链接脚本、BSS 清零、早期栈和内核入口。
\end{enumerate}

可运行模型建议顺序如下：

\begin{enumerate}
\tightlist
\item \filepath{ch00-stage-one-gap-closure.tex}：检查第一阶段还缺哪些能力。
\item \filepath{ch00-stage-one-hardware-os-mapping.tex}：把硬件事实映射成内核对象。
\item \filepath{ch00-stage-one-cpp20-hardware-model.tex} 和 \filepath{ch00-stage-one-x86-64-contract-model.tex}：用可编译模型验证权限、MMIO、中断、DMA 和 x86-64 fault 语义。
\item \filepath{ch00-stage-one-evidence-and-acceptance.tex}：用测试和观测证据验收理解。
\end{enumerate}

主线章节的压缩归并关系如下：

\begin{longtable}{@{}L{0.28\linewidth}L{0.60\linewidth}@{}}
\toprule
主章 & 已压缩吸收的补充方向 \\
\midrule
进程、线程与调度 & 调度器实现、进程生命周期、同步状态机、futex、锁、死锁、调度观测和源码实验。 \\
系统调用、权限与内核边界 & x86-64 入口、用户指针复制、异常修复表、capability、LSM、seccomp、vDSO、审计和资源限制。 \\
exec、ELF 装载与用户栈 & ELF/binfmt、解释器脚本、动态链接器、auxv、CLOEXEC、setuid/capability、ASLR、\code{vfork/posix\_spawn} 和多线程 exec 边界。 \\
虚拟内存与地址空间 & page fault 走读、COW、VMA/PTE、TLB、THP、swap、mlock、PSI、OOM、内核分配器和 C++20 内存模型素材。 \\
mmap、page cache 与文件映射 & \code{do\_mmap}、XArray、writeback、\code{O\_DIRECT}、DAX、映射截断、\code{msync/fsync} 和持久化恢复边界。 \\
设备驱动、中断下半部与 DMA & Linux 设备模型、PCI、virtio、NVMe、IOMMU、DMA 生命周期、块层 request、flush/FUA 和提交顺序。 \\
文件、设备与 I/O & VFS、路径解析、fd/file/inode、page cache、epoll、io\_uring、文件系统日志、目录项持久化和恢复协议。 \\
socket、网络入口与内核边界 & skb、NAPI、Netfilter、socket option、connect/recv 错误语义、小包延迟、超时重试和业务幂等。 \\
安全、隔离与最小权限 & DAC/MAC、capability、namespace、cgroup、Landlock、seccomp、BPF LSM、panic、错误路径、观测、调试和全册验收标准。 \\
\bottomrule
\end{longtable}
